% !TEX root = ../../main.tex
% App H: Transverse-traceless gauge admissibility in the presence of matter.
% Forward-reference targets:
%   \cref{TTGaugeAdmissibility} from theory.tex §2.2 and discussion.tex §4
%   (when the open-questions paragraph is drafted).

\section{Transverse-traceless gauge in the presence of matter}%
\label{TTGaugeAdmissibility}

The standard derivation of the
transverse-traceless~(TT) gauge for the metric perturbation
$\FieldH{_\mu_\nu}$ fails in the presence of matter,
justifying the gauge-free evolution adopted in \cref{Linearisation}.
We follow Hobson~\etal~\cite[Ch.~17--18]{hobson2006general}, with the same result in
Maggiore~\cite[\S1.2]{maggiore2007gravitational} and Misner~\etal~\cite[Box~35.1]{misner1973gravitation}.

\paragraph*{Lorenz gauge}
Writing $g_{\mu\nu} = \eta_{\mu\nu} + \FieldH{_\mu_\nu}$, where $|\FieldH{_\mu_\nu}| \ll 1$
\cite[eq.~(17.1)]{hobson2006general}, an infinitesimal
coordinate transformation $x^{\mu}\to x^{\mu}+\xi^{\mu}(x)$ acts on
the perturbation as
\begin{align}
  h'_{\mu\nu} &= h_{\mu\nu}
    - \partial_{\mu}\xi_{\nu} - \partial_{\nu}\xi_{\mu} ,
  \label{TTGGaugeTransform}
\end{align}
\cite[eq.~(17.4)]{hobson2006general}. With the trace-reversed
perturbation
$\bar{h}_{\mu\nu}\equiv\FieldH{_\mu_\nu}-\tfrac{1}{2}\eta_{\mu\nu}h$,
the gauge transformation
\cref{TTGGaugeTransform} shifts
$\partial^{\mu}\bar{h}_{\mu\nu}\to\partial^{\mu}\bar{h}_{\mu\nu}-\Box^{2}\xi_{\nu}$.
Choosing $\xi^{\mu}$ to satisfy
$\Box^{2}\xi^{\mu}=\partial_{\rho}\bar{h}^{\mu\rho}$ imposes the
Lorenz condition $\partial^{\mu}\bar{h}_{\mu\nu}=0$~\cite[eq. (17.14)]{hobson2006general} while reducing the
linearised Einstein field
equation~\cite[eq.~(17.11)]{hobson2006general} to
\begin{align}
  \Box^{2}\bar{h}^{\mu\nu} &= -2\kappa_{(\mathrm{E})}\,T^{\mu\nu} .
  \label{TTGLorenzWaveEq}
\end{align}

\paragraph*{Failure in the presence of matter}
In vacuum, further residual
gauge transformations preserving $\Box^{2}\xi^{\mu}=0$ set four
linear combinations of $\bar{h}_{\mu\nu}$ to zero, yielding the TT
gauge $\bar{h}_{\mathrm{TT}}^{\,0i}=\bar{h}_{\mathrm{TT}}=0$~\cite[\S18.3,~eq.~(18.8)]{hobson2006general}.
In the presence of matter \cref{TTGLorenzWaveEq} sources an
inhomogeneous solution for $\bar{h}_{\mu\nu}$, and the residual gauge
transformations cannot annul the same combinations of the entire
$\bar{h}_{\mu\nu}$ solution.

For the standard Gertsenshtein configuration of \cref{Linearisation} the
matter is the static external field $\bar{\bm{B}}=\Bzero\hat{x}$,
with
\begin{align*}
  T_{00} &= \tfrac{1}{2}\Bzero^{2} , &
  T_{0i} &= 0 , &
  T_{ij} &= \tfrac{1}{2}\Bzero^{2}\delta_{ij}
           - \bar{B}_{i}\bar{B}_{j} , &
  T^{\mu}{}_{\mu} &= 0 .
  % \label{TTGMaxwellStress}
\end{align*}
The gravitational wave equation
\cref{TTGLorenzWaveEq} for $\bar{h}_{ij}$ acquires a
$h$-dependent source $\delta T_{ij}\supset\tfrac{1}{2}\bar{B}^{2}
h_{ij}$~\cite[eq. (13)]{hwangnoh2023graviton}.
The conventional TT-gauge
result survives this back-reaction only in the limit of a
plane wave on a uniform background.

\paragraph*{Independent failure in $\MAGFt{}^{2}$ theories}
In GR, $\bar{h}_{\mu\nu}$ obeys $\Box^{2}\bar{h}_{\mu\nu}=0$ in vacuum, and the residual
gauge generators $\xi^{\mu}$ satisfy the same second-order equation;
per Fourier mode the four free components of $\xi^{\mu}$ are exactly
the right count to remove four $\bar{h}_{\mu\nu}$ combinations.
A curvature-squared term in the action gives $\bar{h}_{\mu\nu}$ a
higher-derivative equation of motion while $\xi^{\mu}$ stays
second-order; the new equation propagates an additional scalar mode
that the residual gauge choice cannot
remove~\cite{sotiriou2010fr,sezgin1980ghostfree,lin2019pgt}.
Setting \eg~$h^{i}{}_{i}=0$ in such a theory therefore removes a
physical degree of freedom.

An issue raised by
Ejlli~\cite[footnote on p.~2]{ejlli2020graviton} is whether imposing
the gauge conditions in the action and on the EOM give
the same graviton--photon dynamics.
% A Lagrangian-level substitution
% discards the non-radiative components the constraint equations and
% the magnetised stress-energy would otherwise source, so any
% difference between the two routes resides in those components.
\TIDAL{} can derive both and provides a direct empirical resolution;
however, this is beyond the scope of the present work.
